\section*{Bab \ref{ch:g12:exp} --- Eksponen dan Logaritma}

\begin{solution}{exo:g12:exp:1}
$\dfrac{\eu^{3x}\,\eu^{-x+1}}{\eu^{x}} = \eu^{3x - x + 1 - x} = \eu^{x+1}$.

$\ln\!\left(\dfrac{\eu^2\sqrt{\eu}}{\eu^{-3}}\right)
= 2 + \tfrac12 + 3 = \tfrac{11}{2}$.
\end{solution}

\begin{solution}{exo:g12:exp:2}
\emph{(a)} Ambil $X = \eu^x > 0$: $X^2 - 3X + 2 = 0$ memberi $X = 1$ atau
$X = 2$, keduanya positif, sehingga $x = 0$ atau $x = \ln 2$.

\emph{(b)} \omterm{def:g11:func:function}{Daerah asalnya}: $x - 1 > 0$ dan $x + 2 > 0$, jadi $x > 1$.
\omterm{def:g10:algebra:equation}{Persamaannya} menjadi $\ln\bigl((x-1)(x+2)\bigr) = \ln 4$, sehingga
$(x-1)(x+2) = 4$, \emph{yaitu} $x^2 + x - 6 = 0$, jadi $x = 2$ atau $x = -3$.
Hanya $x = 2$ yang berada di \omterm{def:g11:func:function}{daerah asalnya}: $S = \{2\}$.
\end{solution}

\begin{solution}{exo:g12:exp:3}
Bagilah dengan $\eu^x$:
$\dfrac{1 - x^2\eu^{-x}}{1 + \eu^{-x}} \to \dfrac{1-0}{1+0} = 1$, dengan
memakai $x^2 \eu^{-x} \to 0$ (\cref{thm:g12:exp:growth}).

$\dfrac{\ln x}{\sqrt x} \to 0$ menurut \cref{prop:g12:exp:lnanalytic} (kasus
$n = 2$).

$x^2 \ln x = x \cdot (x \ln x) \to 0 \times 0 = 0$.

$x - \ln x = x\left(1 - \frac{\ln x}{x}\right) \to +\infty$ sebab
$\frac{\ln x}{x} \to 0$.
\end{solution}

\begin{solution}{exo:g12:exp:4}
$f'(x) = \eu^{-x} - x\eu^{-x} = (1 - x)\eu^{-x}$, yang bertanda sama dengan
$1 - x$: jadi $f$ naik pada $\intoc{-\infty}{1}$, turun pada
$\intco{1}{+\infty}$, dengan \omterm{def:g10:functions:extrema}{maksimum} global $f(1) = \eu^{-1}$.

Limitnya: ketika $x \to +\infty$, $f(x) = \frac{x}{\eu^x} \to 0$
(\cref{thm:g12:exp:growth}); ketika $x \to -\infty$, $x \to -\infty$ dan
$\eu^{-x} \to +\infty$, sehingga $f(x) \to -\infty$.

$f''(x) = -\eu^{-x} - (1-x)\eu^{-x} = (x - 2)\eu^{-x}$, yang berganti tanda di
$x = 2$: jadi \omterm{def:g12:deriv:inflection}{titik beloknya} $\bigl(2,\, 2\eu^{-2}\bigr)$.
\end{solution}

\begin{solution}{exo:g12:exp:5}
Misalkan $g(x) = x - \ln(1+x)$ pada $\intoo{-1}{+\infty}$. Maka
$g'(x) = 1 - \frac{1}{1+x} = \frac{x}{1+x}$, yang negatif pada
$\intoo{-1}{0}$ dan positif pada $\intoo{0}{+\infty}$: jadi $g$ mencapai
\omterm{def:g10:functions:extrema}{minimumnya} $g(0) = 0$, sehingga $g \geq 0$ dengan kesamaan hanya di $0$. Karena
itu $\ln(1+x) \leq x$.

Terapkan ini dengan $x = \frac1n$:
$n\ln\left(1 + \frac1n\right) \leq 1$, sehingga
$\left(1+\frac1n\right)^n = \eu^{\,n\ln(1+1/n)} \leq \eu$.
Terapkan pula dengan $x = -\frac{1}{n+1} > -1$:
$\ln\left(1 - \frac{1}{n+1}\right) \leq -\frac{1}{n+1}$, sehingga
$-(n+1)\ln\left(1 - \frac{1}{n+1}\right) \geq 1$ dan
$\left(1 - \frac{1}{n+1}\right)^{-(n+1)} \geq \eu$.
\end{solution}

\begin{solution}{exo:g12:exp:6}
\emph{1.} Setiap tahunnya mengalikan modalnya dengan $1.03$: $C_n = C_0 (1.03)^n$.

\emph{2.} $C_n \geq 2C_0 \iff (1.03)^n \geq 2 \iff n \ln 1.03 \geq \ln 2
\iff n \geq \frac{\ln 2}{\ln 1.03} \approx 23.45$: jadi modalnya berlipat dua
setelah $24$ tahun.

\emph{3.} $\frac{70}{3} \approx 23.3$, dekat dengan nilai eksak
$\frac{\ln 2}{\ln 1.03}$. Kaidah itu berhasil sebab
$\ln(1 + r) \approx r$ untuk $r$ yang kecil, sehingga
$\frac{\ln 2}{\ln(1+r)} \approx \frac{0.693}{r} \approx \frac{70}{100r}$.
\end{solution}

\begin{solution}{exo:g12:exp:7}
$\eu^{2x} > \eu^{x+1} \iff 2x > x + 1 \iff x > 1$ (\omterm{thm:g12:exp:existence}{fungsi eksponen} tegas
naik): jadi $S = \intoo{1}{+\infty}$.

\omterm{def:g11:func:function}{Daerah asal} pertidaksamaan kedua: $x^2 - 1 > 0$ dan $x + 5 > 0$, \emph{yaitu}
$x \in \intoo{-5}{-1} \cup \intoo{1}{+\infty}$. Pada daerah itu, karena $\ln$
tegas naik,
\[
\ln(x^2-1) \leq \ln(x+5) \iff x^2 - 1 \leq x + 5 \iff x^2 - x - 6 \leq 0
\iff x \in \intcc{-2}{3}.
\]
Setelah diiriskan dengan \omterm{def:g11:func:function}{daerah asalnya}: $S = \intco{-2}{-1} \cup \intoc{1}{3}$.
\end{solution}

\begin{solution}{exo:g12:exp:8}
\emph{1.} $f'(x) = \dfrac{\eu^x x - \eu^x}{x^2} = \dfrac{(x-1)\eu^x}{x^2}$,
yang negatif pada $\intoo{0}{1}$ dan positif pada $\intoo{1}{+\infty}$: jadi
\omterm{def:g10:functions:extrema}{minimumnya} $f(1) = \eu$. Limitnya: $f \to +\infty$ di $0^+$ (pembilangnya
$\to 1$, penyebutnya $\to 0^+$) dan di $+\infty$ (\cref{thm:g12:exp:growth}).

\emph{2.} Untuk $x > 0$, $\eu^x = kx \iff f(x) = k$. Dari \omterm{def:g10:functions:table}{tabel variasinya}
($+\infty \searrow \eu \nearrow +\infty$, \omterm{def:g12:limcont:continuity}{kontinu} dan tegas \omterm{def:g12:seq:monotonic}{monoton}
pada masing-masing sisinya): tanpa penyelesaian untuk $k < \eu$; tepat satu
($x = 1$) untuk $k = \eu$; dan tepat dua untuk $k > \eu$ (satu di
$\intoo{0}{1}$, satu di $\intoo{1}{+\infty}$, menurut teorema bijeksi pada
masing-masing \omterm{def:g10:numbers:interval}{interval}).
\end{solution}

\begin{solution}{exo:g12:exp:9}
\emph{1.} Langsung dari ketaksamaan pertama pada \cref{exo:g12:exp:5}:
$u_n = \eu^{\,n\ln(1+1/n)} \leq \eu^1 = \eu$.

\emph{2.} Tulislah
\[
\ln u_n = n \ln\!\left(1 + \frac1n\right)
= \frac{\ln\!\left(1 + \frac1n\right) - \ln 1}{\frac1n},
\]
yaitu hasil bagi selisih $\ln$ di titik $1$ dengan pertambahan
$h = \frac1n \to 0$. Karena $\ln$ \omterm{def:g12:deriv:derivative}{dapat diturunkan} di $1$ dengan \omterm{def:g12:deriv:derivative}{turunan}
$1$, maka $\ln u_n \to 1$. Menurut \omterm{def:g12:limcont:continuity}{kekontinuan} $\exp$
(\cref{prop:g12:limcont:seqcont}), $u_n = \eu^{\ln u_n} \to \eu^1 = \eu$.
\end{solution}

\begin{solution}{pb:g12:exp:1}
\textbf{1.} $x = \frac{\ln 5}{2}$. \quad
$3x - 1 = \eu^2$: jadi $x = \frac{\eu^2 + 1}{3}$. \quad Dengan
$u = \eu^x$: $u^2 - 3u + 2 = (u - 1)(u - 2) = 0$: jadi $x = 0$ atau
$x = \ln 2$.

\textbf{2.} $\frac{3\ln 2}{\ln 2} = 3$; \quad
$3 + \frac12 = \frac72$; \quad $\frac52$.

\textbf{3.} $f'(x) = \eu^x - 1$: negatif sebelum $0$, positif
sesudahnya: jadi \omterm{def:g10:functions:extrema}{minimumnya} $f(0) = 0$, sehingga $f \geq 0$ di mana-mana:
$\eu^x \geq 1 + x$. Dengan mensubstitusikan $x = \ln(1 + u)$:
$1 + u \geq 1 + \ln(1 + u)$, yaitu $\ln(1 + u) \leq u$.

\textbf{4.} $x^2\eu^{-x} = \frac{x^2}{\eu^x} \to 0$ dan
$\frac{\ln x}{x} \to 0$: \omterm{thm:g12:exp:existence}{fungsi eksponen} meremukkan pangkat, dan pangkat
meremukkan \omterm{def:g12:exp:ln}{logaritma}. Lalu $\frac{\eu^x - 1}{x} =
\frac{\eu^x - \eu^0}{x - 0} \to \exp'(0) = 1$.

\textbf{5.} $f'(x) = \ln x + 1$: bernilai nol di $\eu^{-1}$, negatif
sebelumnya, positif sesudahnya: jadi \omterm{def:g10:functions:extrema}{minimumnya}
$f\!\left(\frac1\eu\right) = -\frac1\eu$; lalu
$x \ln x \to 0$ di $0^+$: kurvanya meninggalkan titik asal, menukik ke
$-\frac1\eu$, lalu mendaki menjauh.

\textbf{6.} $n = \frac{\ln 2}{\ln 1.03} \approx 23.4$ tahun.

\textbf{7.} Eksaknya: $23.4$, $11.9$, $8.0$ tahun; menurut kaidah 72:
$24$, $12$, $8$. Karena $\ln(1+r) \approx r$, berlaku
$\frac{\ln 2}{\ln(1+r)} \approx \frac{0.693}{r}$, yaitu
$\frac{69.3}{\text{laju dalam }\%}$; para bankir membulatkannya ke atas
menjadi $72$ sebab bilangan itu terbagi dengan indah oleh
$2, 3, 4, 6, 8, 9, 12$ --- berhitung di kepala mengalahkan satu angka
ketelitian.

\textbf{8.} $2^{t/20} = \eu^{t \ln 2 / 20}$:
jadi $\lambda = \frac{\ln 2}{20}$ per menit. Setelah
$24 \times 60 = 1440$ menit: $2^{72} \approx 4.7 \times
10^{21}$ bakteri --- lebih banyak daripada butiran pasir di Bumi,
dari satu sel dalam sehari. Modelnya jujur hanya selama makanan
dan ruangnya bertahan: pertumbuhan yang sesungguhnya membelok menjadi kurva
S logistik (kurva ahli epidemiologi pada \cref{pb:g12:deriv:1}).

\textbf{9.} Pada pukul 23.00 (setelah $8$ jam):
$100 \times 2^{-8/5} \approx 33$ mg. Di bawah $10$ mg:
$2^{-t/5} < 0.1$ memberi $t > 5\,\frac{\ln 10}{\ln 2} \approx
16.6$ jam: jadi sekitar pukul 07.40 keesokan paginya --- espreso pukul
tiga berekor panjang.

\textbf{10.} Setiap tahun: $2$. Setiap bulan:
$\left(1 + \frac{1}{12}\right)^{12} \approx 2.613$. Setiap hari:
$\approx 2.715$. Langit-langitnya: $\eu = 2.71828\ldots$
(\cref{exo:g12:exp:9}) --- dengan pemajemukan sinambung, ketamakannya
\omterm{def:g12:seq:limit}{konvergen}.

\textbf{11.} Jika kasusnya tumbuh secara eksponensial, $c(t) = c_0
\eu^{\lambda t}$, maka $\ln c(t) = \ln c_0 + \lambda t$: yaitu
garis lurus bergradien $\lambda$ --- yaitu laju pertumbuhannya. Bentangan
lurus pada gambar logaritmis itu \emph{adalah} tahap
eksponensialnya, dan kecuramannya adalah tempo wabahnya.

\textbf{12.} $120 - 60 = 60$ dB berarti
$10\log_{10}(I_2/I_1) = 60$: jadi konsernya
$10^6$ --- sejuta kali --- lebih kuat daripada
percakapannya.

\textbf{13.} Amplitudonya: $10^{7-5} = 100$. Energinya:
$100^{3/2} = 1000$: jadi dua angka magnitudo menyembunyikan tiga orde
besaran energi.

\textbf{14.} $10^{6.5 - 2} = 10^{4.5} \approx 32\,000$: jadi air
jeruk lemon tiga puluh ribu kali lebih asam daripada susu --- pH
memampatkan jurang kimia ke dalam skala saku.

\textbf{15.} $\log_2\!\frac{4186}{27.5} =
\frac{\ln(4186/27.5)}{\ln 2} \approx 7.25$: jadi piano membentang sedikit
lebih dari tujuh oktaf. Indra menanggapi \emph{perbandingan} rangsangan
--- melipatduakan intensitasnya terasa seperti satu langkah, berapa pun
taraf awalnya --- sehingga persepsi dibangun atas skala logaritmis,
dan begitu pula satuan yang kita ciptakan untuknya.

\textbf{16.} Menempatkan batang untuk $a$ ujung bertemu ujung dengan
kedudukan batang untuk $b$ menjumlahkan panjangnya $\ln a + \ln b$, dan
tanda yang duduk pada panjang total itu terbaca
$\eu^{\ln a + \ln b} = ab$: jadi mistar hitung menghitung hasil kali
dengan menjumlahkan \omterm{def:g12:exp:ln}{logaritma} --- $\ln(ab) = \ln a + \ln b$
(\cref{prop:g12:exp:lnalg}) yang terpahat pada kayu.

\textbf{17.} \omterm{def:g10:functions:extrema}{Minimum} $\frac{\eu^x}{x}$ pada
$\intoo{0}{+\infty}$ adalah $\eu$, di $x = 1$: jadi tanpa penyelesaian untuk
$k < \eu$, tepat satu untuk $k = \eu$, dan dua untuk $k > \eu$.
Pemeriksaan \omterm{def:g12:deriv:tangent}{garis singgungnya}: di $a = 1$ \omterm{def:g12:deriv:tangent}{garis singgung} $\eu^x$ adalah
$y = \eu + \eu(x - 1) = \eu x$: garis itu melalui titik asal
--- yaitu garis ambangnya, yang menyerempet kurvanya di $(1, \eu)$.

\textbf{18.} $n \ln\!\left(1 - \frac1n\right) =
\frac{\ln(1 - 1/n)}{-1/n} \times (-1) \to -1$, sehingga
$\left(1 - \frac1n\right)^n \to \eu^{-1} \approx 0.368$: jadi
$0.999^{1000} \approx 0.368$ pada \cref{pb:g11:binom:1} pun
terjelaskan --- tetapan nyaris-meleset milik undian itu adalah
$\frac1\eu$.

\textbf{19.} Keduanya adalah bentuk limit ``banyak \omterm{def:g11:prob:model}{kejadian} saling
bebas yang masing-masing tidak mungkin'': peluang bahwa
\emph{tidak satu pun} dari $n$ peluang berukuran $\frac1n$ terwujud menuju
$\eu^{-1}$, dan kaidah sekretaris menyetel jendela penolakannya
sehingga keberhasilannya memusat pada tetapan yang sama itu.
$\frac1\eu \approx 0.368$.

\textbf{20.} Langit-langit pemajemukan; satu-satunya bilangan pokok yang
\omterm{def:g12:deriv:tangent}{garis singgungnya} di $0$ bergradien $1$ (itulah sebabnya kalkulus mencintainya);
pertumbuhan yang melampaui setiap pangkat; tetapan tempat
nyaris-meleset dan penghentian optimum mengendap; serta \omterm{def:g12:exp:ln}{logaritmanya} ---
perkalian yang menjadi penjumlahan, dan rentang paling liar di dunia
yang terlipat ke atas sebuah penggaris.

\end{solution}
